\SourceLeaf{399}%
\[
\begin{array}{ll}
x+n\tfrac12(x'+x''-2x),&y+n\tfrac12(y'+y''-2y)\\
x'+n\tfrac12(x''+x-2x'),&y'+n\tfrac12(y''+y-2y')\\
x''+n\tfrac12(x+x'-2x''),&y''+n\tfrac12(y+y'-2y'')
\end{array}
\]
or, if we set
\[
\begin{aligned}1-n&=\alpha\\\tfrac12n&=\beta\end{aligned}
\]
they are
\[
\begin{array}{ll}
\alpha x+\beta(x'+x''),&\alpha y+\beta(y'+y'')\\
\alpha x'+\beta(x''+x),&\alpha y'+\beta(y''+y)\\
\alpha x''+\beta(x+x'),&\alpha y''+\beta(y+y')
\end{array}
\]
\par
From the points \(C,C',C''\) drop perpendiculars to \(A'A'',A''A,AA'\) and find the positions of the three points where these perpendiculars intersect. Let the coordinates of the intersection of the last two perpendiculars be
\[
\xi,\quad\eta
\]
They can be determined by the following theorem.

If
\[
\begin{array}{ll}
a,&b\\a',&b'\\a'',&b''\\a''',&b'''
\end{array}
\]
are the coordinates of four points, and the straight line through the first and second points is perpendicular to the line through the third and fourth, we have
\begingroup\small
\[
\begin{aligned}
\tfrac{b'-b}{a'-a}&=\text{tangent of the inclination of the first line to the axis of abscissas}\\
\tfrac{b'''-b''}{a'''-a''}&=\text{tangent of the inclination of the second line to the axis of abscissas}
\end{aligned}
\]
\endgroup
and hence, since one inclination is \(90^0\) greater than the other, the product of the two tangents is \(=-1\), so
\[
\tfrac{b'-b}{a'-a}.\tfrac{b'''-b''}{a'''-a''}=-1
\]
In our case, therefore,
